\section{The variational problem on a capped support}

\subsection*{Normalisations}
The sources work in coordinates $t$ in which the criterion reads $kJ/I>1$
(Song and Yue use the slightly stronger $48J/I_0>2001/2000$, see below). We
use the Polymath scaling $u=4t$, in which the criterion reads $kJ/I>4$: every
length is multiplied by $4$ (so $g=4\delta$ and $C_m=4B_m$ for the threshold
$\delta$ and caps $B_m$ of \cite{Sta26}, and the radius $U_{\mathrm{SY}}=521/2000$
of \cite{SY26} becomes $L=1.042$), and every ratio $kJ/I$ is multiplied by
$4$. Values quoted from the sources have been converted, except where we say
otherwise.

\subsection*{Supports}
Let
$\R_k=\{u\in[0,1]^k:\ u_1+\dots+u_k\le1\}$ and write $L\R_k$ for its dilate.
The one-band supports of \cite{Sta26} (\lbl{Literature}: Definitions~1
and~5 and Proposition~1 of \cite{Sta26}) are described by a length $L$, a marginal radius
$\tau$, a threshold $g$ and caps $C_1\le C_2\le\dots$, with $C_0=0$; with
$\mathcal L(u)=\{i:u_i>g\}$ (the \emph{large set}), $m(u)=|\mathcal L(u)|$ and
$s(u)=\sum_{i\in\mathcal L(u)}u_i$,
\[
T=\{u\in L\R_k:\ s(u)\le C_{m(u)}\}.
\]
For a symmetric $F$ supported on $T$,
\[
I_T(F)=\int_TF^2,\qquad
J_T(F)=\int_{\tau\R_{k-1}}\left(\int F(u',u_k)\,du_k\right)^2du',
\]
and the variational criterion is $kJ_T(F)/I_T(F)>4$. (This one-band form of
$J_T$, with marginals over $\tau\R_{k-1}$, is our transcription of the
criterion of \cite[Prop.~1]{Sta26}; see \cite[\S5]{Kim26}.) We write
$I_0,J_0$ for the same functionals without caps, that is with $T$ replaced by
$L\R_k$. In this scaling the parameters are as follows (\lbl{Literature},
converted).
\begin{center}\small
\begin{tabular}{@{}lllll@{}}
\toprule
 & $k$ & $(L,\tau)$ & $g$ & caps $C_1,C_2,\dots$\\
\midrule
\cite{Sta26} & $49$ & $(1.042,\,0.982)$ & $0.112$ & $0.6,\,0.6,\,0.68,\,0.68,\dots$\\
\cite{Kin26}, \cite{SY26} & $48$ & $(1.042,\,0.982)$ & $91/1250$ & $1899/2500$ for all $m$\\
\cite{Cha26}, Tab.~3 & $45$ & $(53/50,\,249/250)$ & $41/625$ & $4B_{1,m}$ (see below)\\
\cite{Cha26}, L.~9.1 & $45$ & $(53/50,\,249/250)$ & $41/625$ & $0.62,\,0.62,\,0.68,\,0.68,\dots$\\
\bottomrule
\end{tabular}
\end{center}
For \cite{Cha26}, Table~3 lists
$B_{1,m}=\frac1{5000}(777,794,875,917,953,983,1016,1042,1063,1081)$ for
$m\le10$ and $B_{1,m}=1081/5000$ for $m\ge11$, and $C_m=4B_{1,m}$; Lemma~9.1
states $C_1=C_2=31/50$ and $C_m=17/25$ for $m\ge3$. All four cap sets satisfy
$g<C_m\le C_{m+1}\le C_m+g$ for $m\ge1$ (checked by hand and in Python; the
case of the Table~3 caps is also \lbl{Lean-checked}, \texttt{charton\_caps\_valid}).
The inequality fails at $m=0$ for every cap set, which is why
Proposition~\ref{prop:fibre} assumes it only for $m\ge1$.

\subsection*{Exact moments}
All integrals reduce to moments over simplices.

\begin{lemma}[\lbl{Proved here}; elementary]\label{lem:egf}
Let $P_j=\sum_iu_i^j$, $X_1,\dots,X_k$ independent exponential variables of
mean $1$, and $E_k(b)=\mathbb E\prod_jP_j(X)^{b_j}$. Then, with
$w(b)=\sum_jjb_j$,
\[
\int_{\R_k}(1-P_1)^A\prod_jP_j^{b_j}\,du=\frac{A!}{(k+A+w(b))!}\,E_k(b),
\qquad
\sum_bE_k(b)\frac{s^b}{b!}=\left(\sum_c w(c)!\,\frac{s^c}{c!}\right)^{k}.
\]
\end{lemma}

\begin{proof}
The Dirichlet integral gives $\int_{\R_k}(1-P_1)^Au^\alpha\,du=A!\,\alpha!/(k+A+|\alpha|)!$
and $\alpha!=\mathbb EX^\alpha$; sum over the monomials of $\prod_jP_j^{b_j}$,
all of degree $w(b)$. The generating function follows from independence.
\end{proof}

\subsection*{The fibres of a capped support}
Let $R=L-|u'|$ on the fibre over $u'\in\mathbb R^{k-1}_{\ge0}$, and call $u'$
\emph{admissible} if $s(u')\le C_{m(u')}$.

\begin{proposition}[\lbl{Proved here}; elementary]\label{prop:fibre}
Let $C_0=0$ and suppose $C_{m+1}\le C_m+g$ for all $m\ge1$. Then the fibre
$\{t\ge0:(u',t)\in T\}$ is empty when $u'$ is not admissible and equals
$[0,\theta(u')]$ otherwise, where
\[
\theta(u')=\min\bigl(R,\ \max(g,\ C_{m(u')+1}-s(u'))\bigr).
\]
Consequently, with $Q(u',x)=\int_0^xP(u',t)\,dt$ for a polynomial $P$,
\[
J_T(P\,1_T)=J_0(P)-\int_{\substack{|u'|\le\tau\\ u'\text{ not adm.}}}Q(u',R)^2
-\int_{\substack{|u'|\le\tau\\ u'\text{ adm.},\ \theta<R}}\bigl(Q(u',R)^2-Q(u',\theta)^2\bigr),
\]
and $I_T(P\,1_T)=I_0(P)-\int_{L\R_k\setminus T}P^2$.
\end{proposition}

\begin{proof}
If $t\le g$ the large set of $(u',t)$ is $\mathcal L(u')$, so $(u',t)\in T$ if and only
if $u'$ is admissible and $t\le R$. If $t>g$ the large set gains the last
coordinate and the condition is $s(u')+t\le C_{m(u')+1}$. If $m(u')=0$ then
$s(u')=0=C_0$, so $u'$ is admissible. If $m(u')\ge1$, the hypothesis
$C_{m+1}\le C_m+g$ with $t>g$ forces $s(u')<C_{m(u')}$. So in either case a
non-admissible $u'$ has an empty fibre, and for admissible $u'$ the fibre is
$[0,\min(R,g)]\cup(g,\min(R,C_{m+1}-s)]=[0,\theta]$. The marginal of $P\,1_T$ is
$Q(u',\theta)$ on admissible $u'$ and $0$ otherwise; subtract from the cap-free
marginal $Q(u',R)$.
\end{proof}

By symmetry each correction is a finite signed sum, over the size $m$ of the
large set and over $j\ge0$ further coordinates forced above $g$ (inclusion and
exclusion for ``all other coordinates are at most $g$''), of integrals of
polynomials over polygons in two radial variables; so every term is an exact
rational number. The ranges that occur are short. In Lean
(\lbl{Lean-checked}, file \texttt{lean/TwinPrimes/Caps.lean}): at most $9$
coordinates of a point of $L\R_{49}$ exceed $g$ for \cite{Sta26}
(\texttt{stadlmann\_large\_count}), at most $10$ large coordinates of a point
of $T$ for \cite{Kin26} (\texttt{kintali\_large\_count}), and at most $16$
coordinates of a point of $L\R_{45}$ exceed $g$ for \cite{Cha26}
(\texttt{charton\_large\_count}); and the inclusion--exclusion over further
coordinates stops at $j\le3$, $3$, $6$ respectively
(\texttt{stadlmann\_extra\_le\_three}, \texttt{kintali\_extra\_le\_three},
\texttt{charton\_extra\_le\_six}). The correction pieces themselves include
non-admissible bases, whose size $m$ is bounded only by the total length (up
to about $14$ for \cite{Kin26}). The reduction to two radial variables is
written out in \cite[\S5]{Kim26}.

